% !TEX program = pdflatex
% Correction — DS Physique-Chimie — 2nde Bac Pro — Chapitre 10
% Température et capteurs thermiques — Filière bois (atelier de menuiserie)

\documentclass[a4paper,11pt]{article}
\usepackage[utf8]{inputenc}
\usepackage[T1]{fontenc}
\usepackage[french]{babel}
\usepackage{geometry}
\geometry{left=2cm,right=2cm,top=1.8cm,bottom=2cm}
\usepackage{amsmath,amssymb}
\usepackage{enumitem}
\usepackage{array}
\usepackage{booktabs}
\usepackage{xcolor}
\usepackage{tcolorbox}
\usepackage{fancyhdr}
\usepackage{lastpage}

\tcbuselibrary{skins,breakable}

% Couleurs (identiques au sujet)
\definecolor{violet}{HTML}{6f42c1}
\definecolor{violetclair}{HTML}{f5f0ff}
\definecolor{orange}{HTML}{d97706}
\definecolor{bleu}{HTML}{0056b3}
\definecolor{vert}{HTML}{2da44e}
\definecolor{gris}{HTML}{718096}

% En-tête / pied
\pagestyle{fancy}
\fancyhf{}
\lhead{\small Correction DS Physique-Chimie — 2nde Bac Pro}
\rhead{\small Ch.\,10 — Température et capteurs}
\cfoot{\thepage\,/\,\pageref{LastPage}}
\renewcommand{\headrulewidth}{0.5pt}

\newtcolorbox{bareme}{
  colback=violetclair,
  colframe=violet,
  leftrule=4pt,
  rightrule=0pt,toprule=0pt,bottomrule=0pt,
  boxsep=4pt,left=8pt,right=6pt
}

\begin{document}

\begin{center}
{\LARGE\bfseries\color{violet} Correction — Devoir Surveillé}\\[4pt]
{\large\bfseries Chapitre 10 — Température et capteurs thermiques}\\[4pt]
{\normalsize\itshape Atelier de menuiserie — Filière bois}\\[4pt]
{\normalsize Seconde Baccalauréat Professionnel — Barème : \textbf{/20}}
\end{center}

\vspace{6pt}
\begin{bareme}
\textbf{Barème global :} Exercice 1 (/6) — Exercice 2 (/9) — Exercice 3 (/5) — \textbf{Total : /20}
\end{bareme}

%===============================================================
% EXERCICE 1
%===============================================================
\vspace{6pt}
\noindent{\large\bfseries\color{violet}Exercice 1 — Du Celsius au Kelvin et inversement} \hfill\textit{\small(/6 pt)}

\vspace{2pt}
\textit{Formules utilisées :} $\ T\,(\text{K}) = \theta\,(^\circ\text{C}) + 273\ $ et $\ \theta\,(^\circ\text{C}) = T\,(\text{K}) - 273$.

\begin{enumerate}[label=\textbf{\arabic*.},leftmargin=*,itemsep=4pt]

\item \textbf{Séchoir à bois à $65\ ^\circ$C.} \hfill\textit{\small(1 pt)}
\[ T = 65 + 273 = \mathbf{338\ \text{K}} \]

\item \textbf{Atelier de menuiserie à $18\ ^\circ$C.} \hfill\textit{\small(1 pt)}
\[ T = 18 + 273 = \mathbf{291\ \text{K}} \]

\item \textbf{Conduit de chaudière à copeaux, $T = 313$\,K.} \hfill\textit{\small(1 pt)}
\[ \theta = 313 - 273 = \mathbf{40\ ^\circ\text{C}} \]

\item \textbf{Local de stockage des colles vinyliques, $T = 268$\,K.} \hfill\textit{\small(1 pt)}
\[ \theta = 268 - 273 = \mathbf{-5\ ^\circ\text{C}} \]
Il s'agit bien d'un stockage au \emph{froid} : les colles vinyliques doivent être maintenues à basse température pour conserver leurs propriétés.

\item \textbf{Tableau de correspondance.} \hfill\textit{\small(1,5 pt)}
\begin{center}
\renewcommand{\arraystretch}{1.35}
\begin{tabular}{|>{\columncolor{violetclair}}c|c|c|c|c|c|}
\hline
\rowcolor{violetclair}
$\theta\ (^\circ$C) & $-273$ & $0$ & $\mathbf{20}$ & $100$ & $\mathbf{300}$ \\
\hline
$T\ $(K) & $\mathbf{0}$ & $\mathbf{273}$ & $293$ & $\mathbf{373}$ & $573$ \\
\hline
\end{tabular}
\end{center}
\textit{Calculs :} $\theta = 293 - 273 = 20\ ^\circ$C ; $T = 0 + 273 = 273$\,K ; $T = 100 + 273 = 373$\,K ; $\theta = 573 - 273 = 300\ ^\circ$C ; $T = -273 + 273 = 0$\,K.

\item \textbf{Impossibilité de $-10$\,K.} \hfill\textit{\small(0,5 pt)}\\
Le \textbf{zéro absolu} vaut $0$\,K : c'est la température la plus basse qui existe dans l'Univers (aucune agitation thermique). \textbf{Aucune température ne peut être négative en Kelvin}, donc $-10$\,K n'a pas de sens physique.

\end{enumerate}

%===============================================================
% EXERCICE 2
%===============================================================
\vspace{10pt}
\noindent{\large\bfseries\color{violet}Exercice 2 — Courbe d'étalonnage de la CTN d'un séchoir à bois} \hfill\textit{\small(/9 pt)}

\begin{enumerate}[label=\textbf{\arabic*.},leftmargin=*,itemsep=4pt]

\item \textbf{Sens de variation et nom du capteur.} \hfill\textit{\small(1,5 pt)}\\
Quand la température $\theta$ augmente (de $20$ à $100\ ^\circ$C), la résistance $R$ \textbf{diminue} (de $10\,000\ \Omega$ à $500\ \Omega$). \\
C'est bien le comportement d'une \textbf{CTN} : \textbf{C}oefficient de \textbf{T}empérature \textbf{N}égatif (la résistance décroît quand la température augmente).

\item \textbf{Lecture : $R = 5\,000\ \Omega$.} \hfill\textit{\small(1,5 pt)}\\
Par lecture directe du tableau ou de la courbe : $R = 5\,000\ \Omega \ \Rightarrow\ \mathbf{\theta = 40\ ^\circ\text{C}}$.

\item \textbf{Consigne $\theta = 80\ ^\circ$C.} \hfill\textit{\small(1,5 pt)}\\
Par lecture du tableau : $\theta = 80\ ^\circ\text{C}\ \Rightarrow\ \mathbf{R = 1\,000\ \Omega}$.\\
Le régulateur du séchoir doit donc commuter le chauffage lorsque la résistance passe sous $1\,000\ \Omega$.

\item \textbf{Conversion de la consigne en Kelvin.} \hfill\textit{\small(1 pt)}
\[ T = 80 + 273 = \mathbf{353\ \text{K}} \]

\item \textbf{Résistance intermédiaire $R = 3\,500\ \Omega$.}

\textbf{a.} \hfill\textit{\small(1 pt)}\\
$3\,500\ \Omega$ est compris entre $R_1 = 5\,000\ \Omega$ (à $40\ ^\circ$C) et $R_2 = 2\,000\ \Omega$ (à $60\ ^\circ$C). \\
La température $\theta$ est donc comprise entre $\mathbf{40\ ^\circ\text{C}}$ et $\mathbf{60\ ^\circ\text{C}}$.

\textbf{b.} Interpolation linéaire. \hfill\textit{\small(2,5 pt)}\\
On applique la formule de l'interpolation linéaire :
\[ \theta = \theta_1 + \frac{R - R_1}{R_2 - R_1}\,(\theta_2 - \theta_1) \]
\[ \theta = 40 + \frac{3\,500 - 5\,000}{2\,000 - 5\,000}\times(60 - 40)
          = 40 + \frac{-1\,500}{-3\,000}\times 20 \]
\[ \theta = 40 + 0{,}5 \times 20 = \mathbf{50\ ^\circ\text{C}}. \]

\end{enumerate}

%===============================================================
% EXERCICE 3
%===============================================================
\vspace{10pt}
\noindent{\large\bfseries\color{violet}Exercice 3 — Thermocouple dans une chaudière à copeaux} \hfill\textit{\small(/5 pt)}

\vspace{2pt}
\textit{Relation donnée :} $\ U\,(\text{mV}) = 0{,}04 \times \theta\,(^\circ\text{C})$.

\begin{enumerate}[label=\textbf{\arabic*.},leftmargin=*,itemsep=4pt]

\item \textbf{Tension à $\theta = 250\ ^\circ$C.} \hfill\textit{\small(1,5 pt)}
\[ U = 0{,}04 \times 250 = \mathbf{10\ \text{mV}} \]

\item \textbf{Température pour $U = 16$\,mV.} \hfill\textit{\small(1,5 pt)}\\
On isole $\theta$ dans la relation :
\[ \theta = \frac{U}{0{,}04} = \frac{16}{0{,}04} = \mathbf{400\ ^\circ\text{C}} \]

\item \textbf{Conversion en Kelvin.} \hfill\textit{\small(1 pt)}
\[ T = 400 + 273 = \mathbf{673\ \text{K}} \]

\item \textbf{Choix du capteur pour $900\ ^\circ$C.} \hfill\textit{\small(1 pt)}\\
Comparaison des plages :
\begin{itemize}[leftmargin=*,itemsep=1pt,topsep=2pt]
  \item \textbf{CTN} : $-50\ ^\circ$C à $+150\ ^\circ$C $\Rightarrow$ \textcolor{gris}{\emph{détruite bien avant $900\ ^\circ$C}}.
  \item \textbf{Pt100} : $-200\ ^\circ$C à $+850\ ^\circ$C $\Rightarrow$ \textcolor{gris}{\emph{en limite de plage, non adaptée}}.
  \item \textbf{Thermocouple type K} : $-200\ ^\circ$C à $+1\,000\ ^\circ$C $\Rightarrow$ \textbf{parfaitement adapté}.
\end{itemize}
\textbf{Conclusion :} on choisit le \textbf{thermocouple type K} pour mesurer la température de la chambre de combustion de la chaudière à copeaux.

\end{enumerate}

\vspace{10pt}
\hrulefill\\[2pt]
\hfill\textit{\small Fin de la correction.}

\end{document}
